\begin{exo}{}
	On considère une fonction $f$ définie sur un intervalle $I$ centré en $0$ et on suppose que $f$ est paire.


	On note $\Cf$ sa courbe représentative dans un repère orthogonal.
	\begin{enumerate}
		\item Rappeler la définition d'une fonction paire.
	\end{enumerate}


	Soient $A$ et $B$, deux points appartenant à $\Cf$ tel que :
	\begin{itemize}
		\item $A$ est le point d'abscisse $x$.
		\item $B$ est le point d'abscisse $-x$.
	\end{itemize}
	\begin{enumerate}[resume]
		\item \begin{enumerate}
			      \item Quelles sont les ordonnées de $A$ et $B$ ?
			      \item Quel lien existe-t-il entre le segment $[AB]$ et l'axe des ordonnées ?
		      \end{enumerate}
		\item Que peut-on en déduire pour $\Cf$ sur $I$ ?
		\item Quelle propriété du cours a-t-on alors démontré ?
	\end{enumerate}
\end{exo}
